% ---
% id: reductions/wpfss/rem:weightfree
% kind: rem
% title: "Why the largest certified weight is used"
% summary: >-
%   Every smaller non-negative weight is also certified, but it gives a cleaner
%   target, which hands the adversary more ([[prop:weight-monotone]]). The
%   largest weight the artefact certifies is therefore the one used, and a
%   smaller one stays valid.
% prominence: standard
% status: accepted
% uses: [def:targets, prop:chain, reductions/cor:monotone, prop:weight-monotone]
% mentions: [req:largest-certified-weight]
% assumes: []
% refs: []
% provenance: newly-derived
% lean:
% ---

\begin{remark}
\label{rem:weightfree}
Nothing in Definition~\ref{def:targets} or Proposition~\ref{prop:chain}
requires $p_v$ to be the largest admissible value: every
$p'_v\in[0,p_v]$ is again a certified weight, and the construction runs
unchanged with $p'=\sum_vp'_v\le p$
(Proposition~\ref{prop:weight-monotone}\,(i)). A smaller weight is not free,
however: it gives a \emph{cleaner} target, which by
Corollary~\ref{cor:monotone} hands the adversary more
(Proposition~\ref{prop:weight-monotone}\,(ii)). The largest weight the
artefact certifies is therefore the best choice, and any smaller one is valid
but weaker (Requirement~\ref{req:largest-certified-weight}).
\end{remark}
